\section{Qwen3 与 Agentic RL：把 reasoning recipe 扩展到真实环境}

前面两个案例仍以数学和代码答案为主，本节把 recipe 扩展到 thinking-mode control 和 agentic tasks。Qwen3 采用多阶段 SFT + reasoning RL，并通过 tags 混合 thinking/non-thinking data，使同一个模型可以控制 test-time compute。

\singleslide{slide-49.jpg}{最终案例 Qwen3：更晚发布的开放 reasoning model，用于观察低数据 RL 与 stage composition。}{49}

Slide 49 不只宣称 Qwen3 排名更高，而是引出两个新的研究问题：reasoning RL 是否必须使用海量题目，以及怎样把窄域推理能力重新并入通用模型。Qwen3 的公开结果提供了一个低数据对照，同时后续 Qwen3-Coder 又把 verifier 扩展到软件工程环境。

\slidepair{slide-50.jpg}{slide-51.jpg}{Qwen3 总体阶段顺序，以及 SFT + reasoning RL 中的 difficulty filtering 与 3,995-example GRPO。}{50--51}

Slide 50 的顺序很关键：reasoning RL 之后还有 general RLHF，最后还能再 distill；它不是单调叠加的单阶段训练。Slide 51 延续 Kimi 的 best-of-$N$ filtering，删除无需 CoT 即可答对的题、清理与验证集相似的问题，并人工检查 CoT 是否真正推导而非猜测。最终 reasoning RL 只用 3,995 个 examples，说明题目质量与在线 rollout 数比静态题目条数更重要。

\begin{importantbox}{“低数据 RL”不等于“低计算 RL”}
3,995 表示 unique prompts 的数量，不表示只生成 3,995 条轨迹。每个 prompt 会被多次采样、验证并随 policy 更新重复使用；真实成本由 group size、平均 CoT 长度、训练轮数和 verifier 开销决定。报告数据效率时必须同时给 rollout tokens 与 compute。
\end{importantbox}

\subsection{Thinking mode fusion}

本节先看推理预算控制。通过特殊 tags 混合短答案与长 reasoning，模型学习何时展开 chain-of-thought；early-stop string 又提供生成长度控制。这相当于把 test-time scaling policy 纳入训练数据协议，而不是让所有任务都无条件生成最长思维链。

\begin{figure}[H]
\centering
\includegraphics[width=0.84\textwidth]{slide-52.jpg}
\caption{Qwen3 thinking-mode fusion：混合 thinking/non-thinking data 并学习提前结束。}
\figsource{CS336 Spring 2026 Lecture 16，Slide 52。}
\end{figure}

\begin{knowledgebox}{术语消化：train-time 与 test-time scaling}
Train-time scaling 增加参数、数据或训练 compute；test-time scaling 让模型在单个问题上生成更长 reasoning、采样更多 candidates 或调用 tools。Thinking-mode fusion 的目标是让模型按任务难度分配 test-time compute。
\end{knowledgebox}

\singleslide{slide-53.jpg}{Test-time scaling：通过更多思考 tokens、更多 samples 或搜索，把推理计算分配给单个问题。}{53}

Slide 53 把 thinking tags 放回更一般的计算视角。对困难题增加 CoT 长度或 candidate 数量，通常能提高 pass@k；但收益会递减，且简单题无条件长思考只增加延迟。理想模型不仅“会长推理”，还应预测什么时候值得继续、什么时候提前停止，并在固定 latency/cost budget 下最大化 verified success。

\subsection{Stage composition 与能力 trade-off}

进一步，课程展示不同 post-training stages 的能力变化：reasoning RL 提升 math/STEM，general RLHF 恢复交互质量，却可能让部分窄域指标回落。Stage order 因而是 multi-objective optimization，不是简单“每加一步都更强”。

\begin{figure}[H]
\centering
\includegraphics[width=0.84\textwidth]{slide-54.jpg}
\caption{不同训练阶段的能力组成：general-purpose alignment 可能轻微牺牲窄域 reasoning。}
\figsource{CS336 Spring 2026 Lecture 16，Slide 54。}
\end{figure}

Slide 54 是 stage composition 的直接证据：reasoning stage 抬高 math/STEM，通用 RLHF 改善聊天、指令与安全后，部分窄域指标略有回落。正确验收不应要求所有指标逐阶段单调增加，而应先定义发布模型的 Pareto frontier：可接受多少 reasoning 退化来换取更好的通用性、简洁度与交互可靠性。

\subsection{Agent environments：从 answer verifier 到 trajectory verifier}

最后把 verifier 从答案扩展到环境轨迹。Qwen3-Coder/agentic RL 使用 repository-level data、pull requests、自动构造的 SWE-bench-style environments 和 tool interaction。Reward 不再只是 final string match，而是 tests、browser state、files changed、commands run 和 task completion。

\singleslide{slide-55.jpg}{Agentic RL 的入口：Qwen3-Coder 在 Qwen3-Next 基础上针对工具使用与软件工程能力后训练。}{55}

数学 RLVR 只需判断最终数值，coding agent 则要在带状态的环境中读取文件、运行命令、修改代码并通过 tests。Trajectory 的每一步都可能改变后续观察，reward 也可能同时包含测试通过、补丁范围、工具错误和效率。Agentic RL 因而更接近传统 sequential decision making，也更容易暴露环境漏洞。

\singleslide{slide-56.jpg}{Qwen3-Coder mid-training 数据：600B repository-level tokens、pull requests、text-code、synthetic QA 与 agent trajectories。}{56}

Slide 56 说明 agent 能力并非从 RL 阶段凭空产生。Repository-level file concatenation 教长上下文代码结构，pull request 配合检索到的仓库状态提供真实修改监督，Common Crawl 的 joint text-code 与 synthetic QA 扩充知识，运行 coding agents 得到的 trajectories 则提前暴露工具协议。RL 主要负责在已有行为支持集上改进决策。

\slidepair{slide-57.jpg}{slide-58.jpg}{从通用 Qwen3-Next 到 Coder 的 expert-model distillation，以及 Web/UX/QA experts 的专项数据。}{57--58}

Slide 57 展示多个 experts：web development、UX、single-turn QA 与 SWE model 分别提供擅长轨迹，再 distill 回统一 Coder policy。Slide 58 解释部分数据质量门槛：web code 由 VLM 与 agent actions 验证可用性，UX expert 学习多种工具格式，QA expert 继续扩大单轮 code synthesis。这里的“expert”既是模型角色，也是数据生成与验证流水线。

\begin{figure}[H]
\centering
\includegraphics[width=0.84\textwidth]{slide-59.jpg}
\caption{自动构造软件工程 agent environments，使训练任务可执行、可验证。}
\figsource{CS336 Spring 2026 Lecture 16，Slide 59。}
\end{figure}

Slide 59 的自动构造路线通常从真实仓库与提交中抽取任务，重建修改前状态，生成问题描述并准备 tests，最终形成类似 SWE-bench 的容器环境。规模可扩展到约 800K tasks，但自动化也会复制 flaky tests、依赖损坏、答案泄漏和任务描述歧义；环境生成器本身需要 held-out audit。

\singleslide{slide-60.jpg}{Agent RL：在可执行环境中生成完整工具轨迹，以测试或任务完成状态提供 reward。}{60}

Agent rollout 的成本由模型 tokens、工具执行、容器启动、依赖安装和测试时间共同决定。训练还必须区分“最终补丁通过”与“轨迹合理”：模型可能读取 git history 找到答案、修改 tests、利用缓存或触发 formal tool 的边界行为。老师明确举例，repository history 与形式验证工具都可能成为 reward-hacking 通道。

\begin{warningbox}{Environment correctness 是 reward correctness}
如果 dependency broken、tests incomplete 或 hidden state 泄漏，agent 会学到错误策略。Agentic RL 的 verifier 需要像 production CI 一样被测试、版本化和监控。
\end{warningbox}

\begin{importantbox}{Agent verifier 的最小防作弊清单}
固定并审计初始 repository state；移除答案提交、隐藏测试与敏感 history；禁止修改 evaluator；隔离网络和跨任务缓存；验证 patch 只触及允许文件；用独立重放确认 reward；对异常短轨迹、测试删除和环境探测行为设置检测。可验证只代表结果可执行检查，不代表检查器对 adversarial policy 安全。
\end{importantbox}

\begin{knowledgebox}{Nota de clase}
课程把 Qwen3-Coder 放在 RLVR 结尾，是为了说明“可验证”正在从数学答案扩展到复杂 environment transitions。越接近真实 agent，reward 越丰富，但环境构造成本和漏洞也越高。
\end{knowledgebox}

\subsection{本章小结}

Qwen3 把 reasoning、thinking control、general alignment 和 agent environments 组合在一起。RLVR 的前沿不只是更好数学题，而是构造可执行、可扩展、能代表真实任务的训练世界。

\singleslide{slide-61.jpg}{全讲回顾：用窄域可验证 reward 缓解偏好过度优化，GRPO 简化训练，但 recipe 与系统仍决定成败。}{61}

最后一页给出的真实结论比“GRPO 很简单”更谨慎。RLHF 与 RLVR 的主要差别是 reward 的可攻击性与成本：可验证 reward 通常更接近目标，但同样可能被漏洞利用；GRPO 已成为需要理解的基础算法，却仍有 bias，在线 RL 也依然 noisy、昂贵且难以调试。成功来自 objective、curriculum、verifier、distillation 与 infra 的共同闭环。

